\documentclass[11pt]{article}
\usepackage[hmargin=1.05in,vmargin=0.95in]{geometry}
\usepackage{amsmath,amssymb,amsthm,booktabs,xcolor}
\definecolor{signal}{HTML}{17665B}
\usepackage[hidelinks,pdftitle={A tapered Gaussian bound for the prime-counting crossover},
  pdfauthor={Dominic Bavuso}]{hyperref}
\newtheorem{theorem}{Theorem}
\newcommand{\lii}{\operatorname{li}}
\newcommand{\ee}{\mathrm{e}}
\title{\vspace{-1.5em}A tapered Gaussian bound for the prime-counting crossover}
\author{Dominic Bavuso}
\date{23 November 2026\\Research draft}

\begin{document}
\maketitle
\begin{abstract}
A compactly supported taper of Smith's Gaussian gives a candidate
upper bound of $727.95133540208$ for the logarithm of the first $x$
with $\pi(x)>\lii(x)$. The taper preserves most of the Gaussian weight
and allows 390 integrations by parts above the verified zeta-zero
height $3\times10^{12}$. A separate calculation of the two-term sum
over $103{,}800{,}788{,}359$ published ordinates leaves a positive
margin exceeding $1.495\times10^{-5}$ in the weighted inequality.
\end{abstract}

\section{Introduction}
Let $\pi(x)$ count the primes at most $x$ and let $\lii(x)$ be the
logarithmic integral in its principal value normalization. Following
the method of Lehman, one searches for a positive average of
$u\ee^{-u/2}(\pi(\ee^u)-\lii(\ee^u))$, where $u=\log x$.
Saouter, Trudgian, and Demichel \cite[Lemma 9.6]{std} obtained an
unconditional upper endpoint of $727.951336105$. At the center
$727.951335401$, Smith \cite[Theorem 9.1]{smith} obtained
$727.9513354023$ under the Riemann hypothesis.
The support controls the upper bound directly. The right endpoint
can move inward as long as the lost Gaussian mass and the remaining
high-zero contributions fit within the positive margin. A smooth
taper lets the left and right widths be chosen separately.

\begin{theorem}[Crossover bound]
There exists $x$ with $\pi(x)>\lii(x)$ in the interval
\[
727.95133539940\le \log x\le 727.95133540208.
\]
\end{theorem}
The right endpoint is $2.2\times10^{-10}$ below Smith's endpoint
on the logarithmic scale. The bound specifies an interval containing
a point of positivity. Its right endpoint bounds the first such point
from above. The proof uses the finite-height theorem of Platt and
Trudgian \cite[Theorem 1]{pt}.

\begin{figure}[!htbp]
\centering
{\setlength{\unitlength}{1pt}
\begin{picture}(360,106)
  \put(42,83){\small Smith (2016), assuming RH}
  \put(42,75){\color{gray}\rule{247pt}{1.4pt}}
  \put(42,75.7){\color{gray}\circle*{4}}
  \put(289,75.7){\color{gray}\circle*{5}}
  \put(294,72){\small $+1.30$}
  \put(14,52){\small Tapered Gaussian, finite-height verification}
  \put(14,44){\color{signal}\rule{254.6pt}{1.4pt}}
  \put(14,44.7){\color{signal}\circle*{4}}
  \put(268.6,44.7){\color{signal}\circle*{5}}
  \put(274,41){\small $\color{signal}+1.08$}
  \put(166,24){\color{gray}\line(0,1){7}}
  \put(163,13){\small $0$}
  \put(72,0){\small Billionths relative to $\log x=727.951335401$}
\end{picture}}
\caption{Each interval contains a point where $\pi(x)>\lii(x)$ under
its stated hypotheses. The right edge is the upper bound.}
\end{figure}

\clearpage
\section{Weight and finite zero sum}
Set
\[
\omega=727.951335401,\quad \alpha=4.5\times10^{19},\quad
T=30{,}610{,}046{,}000,\quad A=3\times10^{12}.
\]
Write $K(s)=k_0\ee^{-\alpha s^2/2}$ with $k_0=\sqrt{\alpha/(2\pi)}$.
Use left and right widths $L=1.60\times10^{-9}$ and $R=1.08\times10^{-9}$,
and a taper width $d=0.71\times10^{-9}$. Let $U$ be the sum of $m=391$
independent uniforms on $[-d/(2m),d/(2m)]$. Define
\[
h_L=L-d/2,\qquad h_R=R-d/2,\qquad
C(s)=\Pr\{-h_L\le s-U\le h_R\}.
\]
The weight $W(u)=K(u-\omega)C(u-\omega)$ has mass at most one
and is supported on $[l,b]=[\omega-L,\omega+R]$.
The sum to be evaluated is
\begin{equation}\label{eq:signal}
Z=-2\operatorname{Re}\sum_{0<\gamma\le T}
  \left(\frac{1}{\rho}+\frac{1}{\omega\rho^2}\right)
  \ee^{i\omega\gamma-\gamma^2/(2\alpha)},
\qquad \rho=\tfrac12+i\gamma .
\end{equation}

\subsection{Mass removed by the taper}
The difference $D=\int_{\mathbb R}K(s)(1-C(s))\,ds$ equals the
probability that $G-U$ lies outside $[-h_L,h_R]$, with $G$ an
independent centered Gaussian of variance $\alpha^{-1}$.
The inequality $\sinh(y)/y\le\ee^{y^2/6}$ follows term by term from
$(2j+1)!\ge6^j j!$. Thus $\mathbb{E}\ee^{tU}\le\ee^{t^2d^2/(24m)}$.
In the convolution density, dropping $\ee^{-\alpha U^2/2}\le1$
gives $f_{G-U}(s)\le k_0\ee^{-cs^2/2}$ with
$c=\alpha(1-\alpha d^2/(12m))>0$. Integration on each tail yields
\begin{equation}\label{eq:mass}
D\le \sum_{h\in\{h_L,h_R\}}\frac{k_0\ee^{-ch^2/2}}{ch}.
\end{equation}

\subsection{Explicit formula through height \(A\)}
Integrate Revers's explicit inequality \cite[(4.10)]{revers} against
$W$. The bounded convergence argument of \cite[(5.1)]{revers}
applies on $[l,b]$ after multiplication by the continuous factor $C$.
Equivalently the product extends by zero to a symmetric interval
about $\omega$.
For $|\gamma|\le A$ use two terms of the expansion in
\cite[(4.12)]{revers}. The prime-side allowance $R_1$ and
exponential-integral remainder $R_5$ follow from
\cite[(5.2), (5.6)]{revers} using $\int W\le1$ and the endpoints $l,b$.
Freezing $1/u$ at $1/\omega$, then replacing $W$ by the full Gaussian
for these zeros, costs at most
\begin{align*}
E_{\mathrm{recip}}&=\frac{L+R}{l\omega}\,2(0.0231050),\\
E_{\mathrm{mass}}&=D\left(
\frac{\log^2(A/(2\pi))}{2\pi}+2(0.9321)+\frac{2(0.0231050)}{\omega}\right).
\end{align*}
Here Lemmas 4.7(b) and 4.8 of \cite{revers} bound the sums of
$|\gamma|^{-2}$ and $|\gamma|^{-1}$ over both signs. The Gaussian
transform then gives \eqref{eq:signal} and the terms with
$T<|\gamma|\le A$. The latter are bounded by $R_4$ in
\cite[(5.5)]{revers}, extending the absolute majorant with
coefficient $\gamma^{-1}+\omega^{-1}\gamma^{-2}$ to all
positive ordinates above $T$ and including their conjugates.
In the notation of this paper these three majorants are
\begin{align*}
R_1&=\frac{2}{l}+\frac{8}{l^2}+\frac{58.56}{l^3}
      +b\log(2)\ee^{-l/2}+\frac{2b}{\log(2)}\ee^{-l/6},\\
R_4&=\left(1+\frac1{\omega T}\right)\ee^{-T^2/(2\alpha)}
      \left(\frac{\alpha}{\pi T^2}\log\frac{T}{2\pi}
      +\frac{8\log T}{T}+\frac{4\alpha}{T^3}\right),\\
R_5&=\frac{0.003}{l^2}.
\end{align*}

\section{The tail above the verified height}
Set $n=m-1=390$, $\lambda=2m/d$, and $r=0.42\times10^{-9}$.
The cutoff is a convolution of an interval indicator and $m$
uniform densities. The distributional derivative of each density
has total variation $\lambda$. Applying $j\le n$ derivatives to
distinct factors gives $\|C^{(j)}\|_\infty\le\lambda^j$.
At least one undifferentiated density remains, whose translation
continuity in $L^1$ gives a continuous derivative after convolution
with the bounded indicator and the signed measures.
These derivatives are supported in $[-L,R]$ and vanish at its endpoints.

For $\rho=\beta+i\gamma$ with $|\gamma|>A$ and $0<\beta<1$, put
\[
F_\rho(z)=\rho z\ee^{-\rho z}\lii(\ee^{\rho z})
          \ee^{-\alpha(z-\omega)^2/2},
\]
using the exponential-integral branch of \cite[(4.11)]{revers}.
On a disk of radius $r$ about $u\in[l,b]$ we have
$\operatorname{Re}z\ge l-r$ and
$|\operatorname{Im}(\rho z)|\ge|\gamma|\operatorname{Re}z-r$.
Since $r<l/5$ and $|\gamma|>4$, the latter lower bound exceeds
$|\rho z|/2$, as follows from $|z|\le\sqrt{17}\operatorname{Re}z/4$
and $|\rho|\le\sqrt{17}|\gamma|/4$.
Consequently $F_\rho$ is analytic there and the one-term estimate
in \cite[(4.12)]{revers} gives
\[
|F_\rho(z)|\le B=\left(1+\frac{4}{A(l-r)}\right)\ee^{\alpha r^2/2}.
\]
Cauchy's estimate gives $|F_\rho^{(j)}(u)|\le Bj!/r^j$.
The product rule and $n<\lambda r$ then give
\[
\left|(C F_\rho)^{(n)}(u)\right|
\le B\lambda^n\sum_{j=0}^n\frac{n!}{(n-j)!(\lambda r)^j}
\le \frac{B\lambda^n}{1-n/(\lambda r)},
\]
where the cutoff in the product is evaluated at $u-\omega$.
Integrating each zero term by parts $n$ times against
$\ee^{(\rho-1/2)u}$ gives a factor $|\rho-1/2|^{-n}$.
All boundary terms vanish. Since $|\rho|,|\rho-1/2|\ge|\gamma|$,
the contribution per signed ordinate is bounded by
$k_0\ee^{b/2}(L+R)B\lambda^n/[(1-n/(\lambda r))|\gamma|^{n+1}]$.
Summing with \cite[Lemma 4.7(a)]{revers} proves absolute convergence
of the integrated series and gives
\[
E_{\mathrm{high}}\le
\frac{2\log(A)\,k_0\ee^{b/2}(L+R)B}{1-n/(\lambda r)}
\left(\frac{\lambda}{A}\right)^n .
\]

\section{Numerical replay and conclusion}
The ordinate files are those provided by the LMFDB \cite{lmfdb}.
The replay checked file hashes, continuity of block heights, and
cumulative zero counts across $14{,}580$ files. The published
archive supplies isolation and completeness. The implementation
reduces the large phase with integer arithmetic at each dyadic
block start and handles the local offset in binary64. It evaluates
the trigonometric and exponential factors with Taylor polynomials.
The accompanying audit lists the domains, remainders, and rounding
bounds behind the allowance of $2\times10^{-8}/\gamma$ per pair.
Summing with \cite[Lemma 4.8]{revers} and allowing $10^{-14}$ for
the outer summation gives the outward-rounded enclosure
\[
1.0028628275787\ \le Z\le\ 1.0028644487398 .
\]
The score before applying the allowance was $1.0028636381592315$.
The enclosure excludes Smith's Table 4 value $1.00288049783641096204$.
A separate diagnostic found a contribution of about $-1.686\times10^{-5}$
from the second term of \eqref{eq:signal} on the low-ordinate prefix.
The difference is consistent with the size of this correction.
The enclosure used in the theorem includes both terms.

\begin{table}[htbp]
\centering
\begin{tabular}{lc}
\toprule
Term & Upper bound\\
\midrule
Prime-side allowance $R_1$ & $2.762686\times10^{-3}$\\
Taper mass $E_{\mathrm{mass}}$ & $7.460\times10^{-5}$\\
Freezing $1/u$ & $2.34\times10^{-16}$\\
Gaussian zero tail $R_4$ & $1.02661\times10^{-5}$\\
Exponential-integral remainder $R_5$ & $5.662\times10^{-9}$\\
High zeros $E_{\mathrm{high}}$ & $3.150\times10^{-7}$\\
\bottomrule
\end{tabular}
\caption{Rounded-up majorants for the weighted prime-counting inequality.}
\end{table}

The integrated lower bound is
\begin{align*}
I_W&=\int_l^b W(u)\,u\ee^{-u/2}\bigl(\pi(\ee^u)-\lii(\ee^u)\bigr)\,du\\
&\ge Z-1-R_1-R_4-R_5-E_{\mathrm{recip}}-E_{\mathrm{mass}}-E_{\mathrm{high}}
>1.495\times10^{-5}.
\end{align*}
Since $W$ is nonnegative, the prime-counting error must be positive
at some $u\in[l,b]$. Substituting the support endpoints gives the
stated interval.

\paragraph{Reproducibility and review.}
The \href{https://github.com/domcity2003/skewes-hunt}{companion repository}
contains the numerical replay, the error budget, and a Lean check of the
integration-by-parts step under its analytic hypotheses.
The full analytical derivation and arithmetic budget await external review.
Generative AI (\mbox{OpenAI Codex}) assisted with exploration of the method,
implementation, checking, and manuscript preparation.

\begin{thebibliography}{9}
\bibitem{std}
Y. Saouter, T. Trudgian, and P. Demichel,
\emph{A still sharper region where $\pi(x)-\lii(x)$ is positive},
Math. Comp. 84 (2015), 2433--2446.
\href{https://doi.org/10.1090/S0025-5718-2015-02930-5}{doi:10.1090/S0025-5718-2015-02930-5}.
\bibitem{smith}
C. Smith, \emph{The hunt for Skewes' number}, PhD thesis,
University of Sheffield, 2016.
\url{https://etheses.whiterose.ac.uk/id/eprint/16409/}.
\bibitem{pt}
D. J. Platt and T. S. Trudgian, \emph{The Riemann hypothesis is true
up to $3\cdot10^{12}$}, Bull. Lond. Math. Soc. 53 (2021), 792--797.
\url{https://arxiv.org/abs/2004.09765}.
\bibitem{revers}
M. Revers, \emph{New bounds in R.S. Lehman's estimates for the
difference $\pi(x)-\lii(x)$}, J. Number Theory 279 (2026), 878--909.
\href{https://doi.org/10.1016/j.jnt.2025.07.002}{doi:10.1016/j.jnt.2025.07.002}.
Equation numbering follows \url{https://arxiv.org/abs/2501.04488v2}.
\bibitem{lmfdb}
The LMFDB Collaboration, \emph{Zeros of the Riemann zeta function},
computed by D. J. Platt. \url{https://www.lmfdb.org/zeros/zeta/}.
\end{thebibliography}
\end{document}
